
\subsubsection{Dataset and Benchmark Catalogs}

Papers With Code aggregates research papers, code repositories, and performance benchmarks across machine learning domains. HuggingFace Datasets Hub extends this infrastructure with programmatic API access and structured metadata schemas. While these catalogs excel at resource discoverability, they lack testability classification logic. This work adds formal verification atop catalog infrastructure via (D,B,M) triple existence checking. The HuggingFace Datasets Hub API provides 100\% metadata completeness (all indexed triples contain required fields) and 84\% coverage of 50 well-known datasets, demonstrating that catalog infrastructure supports automated knowledge base construction without manual curation.

\subsubsection{Confound Documentation}

Documented confounds in deep learning experiments have primarily been treated as domain-specific phenomena. Salesky and Black (2020) identify systematic confounds where tokenizer choice correlates with BLEU scores independent of translation quality improvements. Touvron et al. (2019) reveal confounds between image resolution and model architecture capacity. Goyal et al. (2017) document training procedure confounds, particularly batch size and learning rate coupling. This work demonstrates that confounds reflect fundamental deep learning design choices---preprocessing-metric coupling, architecture-data coupling, hyperparameter interdependencies---that generalize across modalities. Cross-domain transfer validation shows 93.33\% precision across 30 labeled test cases.

\subsubsection{Meta-Research Evaluation Methodologies}

Meta-research tools typically validate performance via expert agreement metrics. Research proposal review systems report 80-85\% inter-rater reliability between expert reviewers. This validation approach suffers from circularity: systems are evaluated against expert consensus rather than ground truth experimental outcomes. This work demonstrates proof-of-concept experimental ground truth validation by testing 20 system-classified ''testable'' hypotheses and measuring post-hoc p-values as success criteria. The 90\% experimental success rate (binomial p = 0.0002 vs. 50\% random baseline) provides non-circular evidence that constraint-satisfiability verification accurately predicts experimental feasibility in proof-of-concept settings.
